% Einheit 2 von 2 – Zinseszins und Guthabentabelle (bank/zinsrechnung/e2.jsonl, zusammenbau v0.1)
%% TODO Zweigzeile Teil 1: was hier gelernt wird (Satz in Schülersprache aus dem Einheitstitel) – Einheit 2
\einheitenkopf[e2]{Einheit 2 von 2 · Zinseszins und Guthabentabelle}
\zweigzeile{ab Kl. 8, je nach Buch bis Kl. 10 (am Gymnasium ab Kl. 7) · P10}
\setcounter{aufgabe}{14}

%% TODO Titel: Ich-kann-Satz fehlt in der Bank (Platzhalter: Kettenname) – Nr. 15
%% TODO Anweisung: ein Satz über der ersten Teilaufgabe fehlt in der Bank – Nr. 15
\begin{aufgabe}{Zinseszins}
\begin{teile}
\teil Auf einem Konto liegen am Anfang 600\,€, nach einem Jahr mit Zinsen 612\,€. Von welchem Betrag werden die Zinsen im zweiten Jahr berechnet? Kreuze an, rechne nicht. \\ \kreuz{vom neuen Guthaben 612\,€} \\ \kreuz{vom Startguthaben 600\,€}
\end{teile}
%% TODO Darstellung für schwach fehlt in der Bank (zinsrechnung-e2-k1-s1-v1; Abschnitt „Für schwache Schüler“)
\swz{900\,€ zu 3\,\%, die Zinsen werden am Jahresende gutgeschrieben – neues Guthaben?}{}
%% TODO Darstellung für schwach fehlt in der Bank (zinsrechnung-e2-k1-s1-v2; Abschnitt „Für schwache Schüler“)
\swz{1\,600\,€ zu 2\,\%, die Zinsen werden am Jahresende gutgeschrieben – neues Guthaben?}{}
%% TODO Darstellung für schwach fehlt in der Bank (zinsrechnung-e2-k1-s1-v3; Abschnitt „Für schwache Schüler“)
\swz{2\,200\,€ zu 5\,\%, die Zinsen werden am Jahresende gutgeschrieben – neues Guthaben?}{}
%% TODO Darstellung für schwach fehlt in der Bank (zinsrechnung-e2-k1-s1-v4; Abschnitt „Für schwache Schüler“)
\swz{350\,€ zu 4\,\%, die Zinsen werden am Jahresende gutgeschrieben – neues Guthaben?}{}
%% TODO Darstellung für schwach fehlt in der Bank (zinsrechnung-e2-k1-s1-v5; Abschnitt „Für schwache Schüler“)
\swz{5\,000\,€ zu 1\,\%, die Zinsen werden am Jahresende gutgeschrieben – neues Guthaben?}{}
\swfrage{Was bleibt gleich, was ändert sich?}
%% TODO Darstellung für schwach fehlt in der Bank (zinsrechnung-e2-k1-s2-v1; Abschnitt „Für schwache Schüler“)
\swz{2\,000\,€ zu 3\,\%, die Zinsen bleiben auf dem Konto. Rechne Jahr für Jahr – Guthaben nach zwei Jahren?}{}
\swa{Zinssatz 2\,\%. Zinsen vom Guthaben der Zeile davor, neues Guthaben = altes Guthaben + Zinsen + Einzahlung. Ergänze Jahr 2. Zinsen \leerfeld[€] , Guthaben \leerfeld[€]}{\sachtabelle{lrrr}{Jahr & Zinsen & Einzahlung & Guthaben}{1 & – & 300,00\,€ & 300,00\,€\\ 2 & \leerzelle & 300,00\,€ & \leerzelle}}
\swa{Zinssatz 2\,\%, jedes Jahr 120\,€ Einzahlung. Ergänze die zwei Felder und kontrolliere mit Jahr 3. Guthaben \leerfeld[€] , Zinsen \leerfeld[€]}{\sachtabelle{lrrr}{Jahr & Zinsen & Einzahlung & Guthaben}{1 & – & – & 350,00\,€\\ 2 & 7,00\,€ & 120,00\,€ & \leerzelle\\ 3 & \leerzelle & 120,00\,€ & 606,54\,€}}
%% TODO Darstellung für schwach fehlt in der Bank (zinsrechnung-e2-k1-s5-v1; Abschnitt „Für schwache Schüler“)
\swz{Rechne in einem Schritt mit dem Wachstumsfaktor: 750\,€ zu 4\,\% – Guthaben nach einem Jahr?}{}
%% TODO Darstellung für schwach fehlt in der Bank (zinsrechnung-e2-k1-s6-v1; Abschnitt „Für schwache Schüler“)
\swz{3\,500\,€ zu 3\,\% mit Zinseszins, 4 Jahre – Guthaben am Ende? Runde auf Cent.}{}
%% TODO Darstellung für schwach fehlt in der Bank (zinsrechnung-e2-k1-s7-v1; Abschnitt „Für schwache Schüler“)
\swz{4\,000\,€ zu 3\,\%, 3 Jahre: einmal mit einfachen Zinsen, einmal mit Zinseszins. Wie viel Euro Unterschied?}{}
%% TODO Darstellung für schwach fehlt in der Bank (zinsrechnung-e2-k1-s8-v1; Abschnitt „Für schwache Schüler“)
\swz{3\,000\,€ zu 2\,\% mit Zinseszins, 5 Jahre – wie viel Euro Zinsen insgesamt?}{}
%% TODO Darstellung für schwach fehlt in der Bank (zinsrechnung-e2-k1-s9-v1; Abschnitt „Für schwache Schüler“)
\swz{Nach 3 Jahren zu 2\,\% mit Zinseszins stehen 5\,306,04\,€ auf dem Konto. Wie viel Euro wurden am Anfang angelegt?}{}
\swa{Für den Führerschein zahlt Lukas jedes Jahr 350\,€ auf ein Konto ein, der Zinssatz ist 3\,\%. Vervollständige die Tabelle und kontrolliere mit Jahr 4. \hfill (P10 2014 OS) Guthaben \leerfeld[€] , Zinsen \leerfeld[€]}{\sachtabelle{lrrr}{Jahr & Zinsen & Einzahlung & Guthaben}{1 & – & 350,00\,€ & 350,00\,€\\ 2 & 10,50\,€ & 350,00\,€ & 710,50\,€\\ 3 & 21,32\,€ & 350,00\,€ & \leerzelle\\ 4 & \leerzelle & 350,00\,€ & 1\,464,27\,€}}
\end{aufgabe}

%% TODO Titel: Ich-kann-Satz fehlt in der Bank (Platzhalter: Kettenname) – Nr. 16
%% TODO Anweisung: ein Satz über der ersten Teilaufgabe fehlt in der Bank – Nr. 16
\begin{aufgabe}{Guthaben nach n Jahren mit jährlicher Einzahlung (Tabelle statt Formel)}
\swa{Start 600\,€, Zinssatz 2\,\%. Am Ende jedes Jahres kommen erst die Zinsen, dann 120\,€ Einzahlung dazu. Fülle die Tabelle aus – Guthaben nach 3 Jahren? \leerfeld[€]}{\sachtabelle{lrrr}{Jahr & Zinsen & Einzahlung & Guthaben}{Start & – & – & 600,00\,€\\ 1 & \leerzelle & 120,00\,€ & \leerzelle\\ 2 & \leerzelle & 120,00\,€ & \leerzelle\\ 3 & \leerzelle & 120,00\,€ & \leerzelle}}
\end{aufgabe}

%% TODO Titel: Ich-kann-Satz fehlt in der Bank (Platzhalter: Kettenname) – Nr. 17
%% TODO Anweisung: ein Satz über der ersten Teilaufgabe fehlt in der Bank – Nr. 17
\begin{aufgabe}{Ratenkauf und Kredit mit Zinseszins (Vorrat, Verbrauchersicht)}
%% TODO Darstellung für schwach fehlt in der Bank (zinsrechnung-e2-k3-s1-v1; Abschnitt „Für schwache Schüler“)
\swz{Ein Kredit über 5\,000\,€ wird nach 3 Jahren auf einmal zurückgezahlt, 6\,\% Zinsen mit Zinseszins. Wie viel Euro sind zurückzuzahlen?}{}
\end{aufgabe}

%% TODO Titel: Ich-kann-Satz fehlt in der Bank (Platzhalter: Kettenname) – Nr. 18
%% TODO Anweisung: ein Satz über der ersten Teilaufgabe fehlt in der Bank – Nr. 18
\begin{aufgabe}{Zinseszins – Fehler finden, Begründen, Anwendung, Darstellungswechsel}
\swz[3]{2\,400\,€ zu 2\,\% mit Zinseszins, 3 Jahre – Guthaben? Jonas rechnet: \rechnung{2\,400 \cdot 0{,}02 &= 48 \\ 2\,400 + 3 \cdot 48 &= 2\,544} Antwort: 2\,544\,€. Finde den Fehler.}{}
\swfrage{Ein Guthaben wächst jedes Jahr um 3\,\%. Begründe, warum es jedes Jahr um mehr Euro wächst.}
\swz[3]{Lara bekommt zur Jugendweihe 2\,000\,€ und legt sie zu 3\,\% mit Zinseszins an. Der Führerschein kostet in 4 Jahren etwa 2\,250\,€. Reicht das Geld?}{}
\swa{Die Tabelle zeigt ein Guthaben mit Zinseszins. Welcher Zinssatz? \leerfeld[\%]}{\sachtabelle{lr}{Jahr & Guthaben}{0 & 2\,000,00\,€\\ 1 & 2\,080,00\,€\\ 2 & 2\,163,20\,€}}
\end{aufgabe}

\uebersichtskasten{Zinseszins: Die Zinsen kommen am Jahresende zum Guthaben dazu und werden im nächsten Jahr mitverzinst – das Guthaben wächst jedes Jahr um denselben Prozentsatz, aber um mehr Euro: 800 € zu 5 \% → nach einem Jahr 840 €, nach zwei Jahren 882 € (nicht 880 €). \\ Wachstumsfaktor: „plus 5 \%“ ist dasselbe wie „mal 1,05“ – q = 1 + p/100: 800 · 1,05 = 840. \\ Endkapital: Startkapital mal Wachstumsfaktor hoch Anzahl der Jahre – die Potenz rechnet der Taschenrechner: $K_n$ = $K_0$ · qⁿ; 800 · 1,05⁶ ≈ 1072,08 €. \\ Tabelle: Zinsen = Zinssatz vom Guthaben der Vorzeile; neues Guthaben = altes Guthaben + Zinsen (+ Einzahlung): 840 € → Zinsen 42 € → 882 €.}
